\documentclass[12pt,a4paper]{article}
\usepackage[utf8]{inputenc}
\usepackage{amsmath,amssymb}
\usepackage{graphicx}
\usepackage[left=1.5cm,right=1cm,top=1.3cm,bottom=1.3cm]{geometry}
\usepackage{lastpage}
% Font Times New Roman (cần XeLaTeX + fontspec — uncomment nếu VPS hỗ trợ):
% \usepackage{fontspec}
% \setmainfont{Times New Roman}

\begin{document}

\begin{center}
\begin{tabular}{|p{0.4\linewidth}|p{0.35\linewidth}|p{0.25\linewidth}|}
\hline
\multicolumn{1}{|c|}{SỞ GDĐT ...} & \multicolumn{2}{c|}{ĐỀ KIỂM TRA} \\
\multicolumn{1}{|c|}{TRƯỜNG THPT ...} & \multicolumn{2}{c|}{MÔN TOÁN 12} \\
\multicolumn{1}{|c|}{Đề chính thức} & \multicolumn{2}{c|}{THỜI GIAN: 90 PHÚT, (không kể thời gian phát đề)} \\
\multicolumn{1}{|c|}{\textit{(Đề thi gồm có .... trang)}} & \multicolumn{2}{c|}{} \\
\hline
\multicolumn{2}{|p{0.75\linewidth}|}{Họ và tên: ........................................................................................ \newline Số báo danh: ....................................................................................} & \multicolumn{1}{c|}{\textbf{Mã đề: 1761}} \\
\hline
\end{tabular}
\end{center}

\bigskip

\begin{center}
{\Large\textbf{ĐÁP ÁN VÀ LỜI GIẢI}}
\end{center}

\noindent\textbf{PHẦN I. Câu trắc nghiệm nhiều phương án lựa chọn.}

\medskip

\noindent\textcolor{blue}{\textbf{Câu 1.}}

\noindent \textit{Đáp án: \textbf{D}.}

\noindent\textbf{Lời giải.}

Từ bảng biến thiên suy ra hàm số đồng biến trên khoảng $(-1;0)$.

\medskip
\noindent\textcolor{blue}{\textbf{Câu 2.}}

\noindent \textit{Đáp án: \textbf{B}.}

\noindent\textbf{Lời giải.}

Dựa vào bảng biến thiên, hàm số đã cho nghịch biến trên các khoảng $(-\infty ;-1)$ và $(0 ; 1)$.

\medskip
\noindent\textbf{PHẦN II. Câu trắc nghiệm đúng sai.}

\medskip

\noindent\textcolor{blue}{\textbf{Câu 1.}}

\noindent \textit{Đáp án: ĐĐSĐ}

\noindent\textbf{Lời giải.}

a)  %Đúng.\\

Dựa vào đồ thị hàm số, ta nhận thấy hàm số đồng biến trên các khoảng $(-\infty;-3)$ và $\left(-1;+\infty\right)$.			
			
b)  %Đúng.\\

			Hàm số đồng biến trên các khoảng  $\left(-1;+\infty\right)$.\\

Mà $2024$; $2025 \in \left(-1;+\infty\right)$  nên $f(2024)<f(2025)$.			
			
c)  %Sai.\\

Đồ thị hàm số có tiệm cận đứng là $x=-2$.\\

Gọi đường tiệm cận xiên có dạng $y=a'x+b'$.\\

Dựa vào đồ thị, ta thấy đường tiệm cận xiên đi qua hai điểm $(0;1)$ và $(-2;-1)$.\\

Suy ra $\begin{cases}b'=1\\-2a'+b'=-1\end{cases} \Leftrightarrow \begin{cases}b'=1\\a'=1.\end{cases}$\\

Vậy đường tiệm cận xiên là $y=x+1$.

d)  %Đúng.\\

	Vì đồ thị hàm số có tiệm cận xiên là $y=x+1$ và tiệm cận đứng là $x=-2$ nên hàm số có dạng $y=x+1+\dfrac{k}{x+2}$.\\

	Vì đồ thị hàm số đi qua điểm $(-1;1)$ nên ta có $1=-1+1+\dfrac{k}{-1+2}\Rightarrow k=1$.\\

	Do đó hàm số đã cho là $y=x+1+\dfrac{1}{x+2}=\dfrac{x^2+3x+3}{x+2}$.\\

	Thay $x=2$ ta tính được $y=\dfrac{13}{4}$ nên đồ thị hàm số đã cho đi qua điểm $A\left(2;\dfrac{13}{4}\right)$.

\medskip
\noindent\textcolor{blue}{\textbf{Câu 2.}}

\noindent \textit{Đáp án: ĐSĐS}

\noindent\textbf{Lời giải.}

Xét $h(t)=6t^3 -81t^2 +324t$ ($0\leq t\leq 8$).\\

		Ta có $h'(t) = 18t^2 -162t +324$; $h'(t) = 0 \Leftrightarrow \left[\begin{array}{l}t=3\\t=6.\end{array}\right.$\\

					Bảng biến thiên $h(t)$
% [Hình TikZ không compile được: tikz_0f209270001b]

		a)  Ta có $h(0) = 0$. Suy ra tại thời điểm ban đầu, quả bóng đang ở mặt đất.
			
b)  Ta có $h(1) = 6\cdot 1 -81\cdot 1+324\cdot 1 = 249$.\\

			Suy ra sau $1$ phút, quả bóng bay đạt độ cao $249$ mét so với mặt đất.
			
c)  Từ bảng biến thiên, ta thấy trong $3$ phút đầu tiên, quả bóng bay tăng dần độ cao.
			
d)  Từ bảng biến thiên ta thấy từ phút thứ $3$ đến phút thứ $6$ quả bóng bay giảm dần độ cao. Sau đó từ phút thứ $6$ đến phút thứ $8$, quả bóng tăng dần độ cao.

\medskip
\noindent\textbf{PHẦN III. Câu trả lời ngắn.}

\medskip

\noindent\textcolor{blue}{\textbf{Câu 1.}}

\noindent \textit{Đáp số: $-3$}

\noindent\textbf{Lời giải.}

Tập xác định: $\mathbb{R}\setminus\left\{-1\right\}$.\\

			Đạo hàm $y'=\dfrac{x^2 + 2 x}{\left(x + 1\right)^2}$;\\
$y'=0 \Leftrightarrow x=-2$ hoặc $x=0$.\\

% [Hình TikZ không compile được: tikz_baaea164774f]

	Hàm số nghịch biến trên $(-2;-1)$ và  $(-1;0)$ nên $a=-2$, $b=-1$ và $c=0$ nên\break $T=a+b+c=-3$.

\medskip
\noindent\textcolor{blue}{\textbf{Câu 2.}}

\noindent \textit{Đáp số: $4$}

\noindent\textbf{Lời giải.}

Ta có $V'=-0,06426+0,0170086T-0,0002307T^2$, $V'=0\Leftrightarrow \left[\begin{array}{l}T\approx 4\in [0;30]\\T\approx 70.\end{array}\right.$\\

		Bảng biến thiên của $V$ trên $[0;30]$.
		
% [Hình TikZ không compile được: tikz_82726dc32fe2]

		Vậy thể tích $V$ tăng trên khoảng $(4;30)$.\\

		Suy ra $T_0=4$.

\medskip
\noindent\textbf{PHẦN IV. Câu tự luận.}

\medskip

\noindent\textcolor{blue}{\textbf{Câu 1.}}

\noindent\textbf{Lời giải.}

Vận tốc của vật là $v(t)=s'(t)=-3t^2+6t+6$.\\

		Ta có $v'(t)=-6t+6$.\\

		Cho $v'(t)=0\Leftrightarrow t=1$.\\

		Bảng biến thiên\\

% [Hình TikZ không compile được: tikz_bcb7f028b670]

		Dựa vào bảng biến thiên, ta thấy $v(t)$ tăng khi $t\in (0;1)$. Vậy vật tăng tốc trong khoảng thời gian từ gian $1$ giây tính từ lúc bắt đầu chuyển động.

\medskip
\noindent\textcolor{blue}{\textbf{Câu 2.}}

\noindent\textbf{Lời giải.}

Tốc độ truyền bệnh $y=f'(t)=-3t^2+90t$ với $t\in [0; 25]$.\\

		Ta có $y'=-6t+90$; $y'=0\Leftrightarrow t=15$.\\

		Bảng biến thiên
		
% [Hình TikZ không compile được: tikz_7d2da6ddaef5]

		Tốc độ truyền bệnh giảm từ ngày $15$ đến ngày $25$.\\

		Do đó $m=15$, $n=25$.\\
Suy ra $m-n=10$.

\medskip

\end{document}